\reading{MR-18}{Fundamental Review of the Trading Book}{strike}{8}
% ==========================================================

\begin{mrkeybox}[title={Learning objectives in this reading}]
\small
\begin{enumerate}[label=\textbf{\alph*.},leftmargin=1.7em]
\item Describe the changes to the Basel framework for calculating market risk capital under the
      Fundamental Review of the Trading Book (FRTB) and the motivations for these changes.
\item Compare the various liquidity horizons proposed by the FRTB for different asset classes
      and explain how a bank can calculate its expected shortfall using the various horizons.
\item Explain the FRTB revisions to Basel regulations in the following areas: classification of
      positions in the trading book compared to the banking book; backtesting; profit and loss
      attribution; credit risk; and securitizations.
\end{enumerate}
\end{mrkeybox}

% ---------------------------------------------------------
\lo{MR-18 a}{Describe the changes to the Basel framework for calculating market risk capital under the Fundamental Review of the Trading Book (FRTB) and the motivations for these changes.}

In May 2012 the Basel Committee on Banking Supervision introduced a consultative document titled
the Fundamental Review of the Trading Book. The final version, published as the \term{Minimum
Capital Requirements for Market Risk}, covered the shift from VaR to expected shortfall, new
methods for calculating capital under stressed conditions, and liquidity horizons varying by
asset class.

\subsection{How the framework evolved}

\begin{center}
\begin{tikzpicture}[every node/.style={font=\scriptsize\sffamily},
  era/.style={rectangle, rounded corners=3pt, line width=0.9pt,
    text width=4.1cm, align=center, inner sep=6pt, minimum height=2.1cm}]
  \node[era, draw=black!45, fill=FRMSoft] (b1) at (0,0)
    {\textbf{Basel I}\\[3pt] Market risk capital based on \textbf{VaR}\\ 10-day horizon\\
     99\% confidence};
  \node[era, draw=FRMGold!70, fill=FRMPaper] (b25) at (5.1,0)
    {\textbf{Basel II.5}\\[3pt] Added a \textbf{stressed VaR} measure\\ over 250 days,
     \textit{in addition to} the current measure};
  \node[era, draw=FRMNavy!60, fill=PaleNavy] (frtb) at (10.2,0)
    {\textbf{FRTB}\\[3pt] \textbf{Stressed expected shortfall}\\ 97.5\% confidence\\
     over a 250-day stressed period};
  \draw[FRMGold, -{Stealth[length=5pt]}, line width=1.1pt] (b1) -- (b25);
  \draw[FRMGold, -{Stealth[length=5pt]}, line width=1.1pt] (b25) -- (frtb);
  \node[anchor=north, text=black!60, align=center, text width=3.4cm] at (2.55,-1.25)
    {\textit{recent losses only}};
  \node[anchor=north, text=black!60, align=center, text width=3.6cm] at (7.65,-1.25)
    {\textit{adds stressed conditions, but still asks ``how bad?''}};
\end{tikzpicture}
\end{center}
\figcap{Two changes happen at the FRTB step, and they are easy to conflate. The \textit{measure}
changes from VaR to ES, and the \textit{confidence level} changes from 99\% to 97.5\%. The
stressed-period requirement carries over, so the FRTB measure is a stressed ES, not a current
one.}

\begin{mrtrapbox}[title={The motivation in one line}]
VaR only asks \term{``how bad can things get?''} --- it reports the threshold and says nothing
about what lies beyond it. ES asks how bad things are \textit{on average} once that threshold is
breached.
\end{mrtrapbox}

\subsection{Problem with VaR, scenario 1: a normal distribution}

For normal distributions with mean $\mu$ and standard deviation $\sigma$, the two measures are
\term{nearly identical}:

\begin{center}
\small
\begin{tabular}{l l}
\arrayrulecolor{FRMRule}\toprule
\rowcolor{FRMNavy} \hdr{Measure} & \hdr{Value}\\
99\% VaR & $\mu + 2.326\,\sigma$\\
\rowcolor{FRMSoft} 97.5\% Expected Shortfall & $\mu + 2.338\,\sigma$\\
\bottomrule
\end{tabular}
\end{center}

\begin{center}
\begin{tikzpicture}
\begin{axis}[name=nrm, width=6.6cm, height=4.6cm,
  xmin=-1, xmax=4.2, ymin=0, ymax=0.45,
  axis y line=none, axis x line=bottom, axis line style={draw=black!55},
  xlabel={\scriptsize Loss, in $\sigma$}, xlabel style={font=\scriptsize},
  xtick={0,1,2,3,4}, tick label style={font=\tiny},
  title={\scriptsize\sffamily\bfseries\color{FRMNavy}Normal: the two coincide},
  clip=false]
\addplot[black!80, line width=1.1pt, smooth, domain=-1:4.2, samples=160]
  {0.39894*exp(-(x^2)/2)};
\addplot[domain=2.326:4.2, samples=60, smooth, draw=none, fill=FRMRed!25, forget plot]
  {0.39894*exp(-(x^2)/2)} \closedcycle;
\draw[FRMNavy, line width=1.0pt] (axis cs:2.326,0) -- (axis cs:2.326,0.095);
\draw[FRMRed, line width=1.0pt] (axis cs:2.338,0) -- (axis cs:2.338,0.075);
\node[anchor=west, font=\tiny\sffamily, text=FRMNavy, fill=white, inner sep=1.2pt]
  at (axis cs:2.5,0.16) {99\% VaR $= 2.326\sigma$};
\node[anchor=west, font=\tiny\sffamily, text=FRMRed, fill=white, inner sep=1.2pt]
  at (axis cs:2.5,0.10) {97.5\% ES $= 2.338\sigma$};
\end{axis}
\begin{axis}[name=fat, at={($(nrm.east)+(1.5cm,0)$)}, anchor=west,
  width=6.6cm, height=4.6cm,
  xmin=-1, xmax=6.5, ymin=0, ymax=0.45,
  axis y line=none, axis x line=bottom, axis line style={draw=black!55},
  xlabel={\scriptsize Loss, in $\sigma$}, xlabel style={font=\scriptsize},
  xtick={0,2,4,6}, tick label style={font=\tiny},
  title={\scriptsize\sffamily\bfseries\color{FRMRed}Fat-tailed: they separate},
  clip=false]
\addplot[black!80, line width=1.1pt, smooth, domain=-1:6.5, samples=200]
  {0.355*exp(-(x^2)/2) + 0.030*exp(-(x^2)/40)};
\addplot[domain=2.20:6.5, samples=90, smooth, draw=none, fill=FRMRed!25, forget plot]
  {0.355*exp(-(x^2)/2) + 0.030*exp(-(x^2)/40)} \closedcycle;
\draw[FRMNavy, line width=1.0pt] (axis cs:2.20,0) -- (axis cs:2.20,0.11);
\draw[FRMRed, line width=1.0pt] (axis cs:4.20,0) -- (axis cs:4.20,0.09);
\node[anchor=west, font=\tiny\sffamily, text=FRMNavy, fill=white, inner sep=1.2pt]
  at (axis cs:2.35,0.20) {99\% VaR};
\node[anchor=west, font=\tiny\sffamily, text=FRMRed, fill=white, inner sep=1.2pt]
  at (axis cs:3.1,0.13) {97.5\% ES --- much further out};
\draw[FRMRed, {Stealth[length=4pt]}-{Stealth[length=4pt]}, line width=0.7pt]
  (axis cs:2.25,0.045) -- (axis cs:4.15,0.045);
\end{axis}
\end{tikzpicture}
\end{center}
\figcap{The shaded region is the tail ES averages over. On the left, a normal distribution: the
tail is thin, so averaging over it lands almost exactly where the VaR cut is, and the two
measures differ by 0.012$\sigma$. On the right, heavier tails: the same tail region now contains
far larger losses, so the average is pulled well beyond the VaR threshold. \textbf{This is the
whole argument for the switch}, and it is invisible in the normal case.}

\subsection{Problem with VaR, scenario 2: a bond portfolio}

\begin{mrexbox}[title={Example --- VaR reports nothing}]
A \$950 million bond portfolio. If it defaults, the loss is 40\% of principal, or
$0.40 \times \$950\text{m} = \mathbf{\$380 \text{ million}}$. The probability of default is
2\%.
\tcblower
\small
\textbf{VaR at 95\%:} the 95th percentile outcome is \textit{no default}, because only 2\% of
outcomes involve a loss at all. So the 95\% VaR is \textbf{\$0} --- VaR reports that the
portfolio is riskless.

\medskip
\textbf{ES:} expected shortfall looks at what happens \textit{beyond} the threshold, and there
the \$380 million default loss is sitting. It is the only thing in the tail, and VaR does not
see it.
\end{mrexbox}

\begin{mrtrapbox}[title={A precision point on this example}]
The source notes state the expected shortfall as \$380 million. Strictly, \$380 million is the
loss \textit{conditional on default}. Expected shortfall is the probability-weighted average
across the whole tail, so with a 2\% default probability inside a 5\% tail it would be
\[ \frac{0.02}{0.05} \times \$380\text{m} \;=\; \$152 \text{ million} \]
The teaching point is unaffected and is the one to carry into the exam: \term{VaR reports \$0
while the portfolio can lose \$380 million}. But know which number is which if a question asks
you to compute ES rather than describe it.
\end{mrtrapbox}

% ---------------------------------------------------------
\lo{MR-18 b}{Compare the various liquidity horizons proposed by the FRTB for different asset classes and explain how a bank can calculate its expected shortfall using the various horizons.}

\begin{mrdefbox}[title={Liquidity horizons}]
FRTB uses \term{five} liquidity horizons: 10, 20, 60, 120 and 250 days. These reflect the time
needed to sell an asset or hedge its risk in \textit{stressed} markets. A 60-day horizon assumes
three months to recover from price volatility.
\end{mrdefbox}

\begin{center}
\begin{tikzpicture}
\begin{axis}[
  width=10.8cm, height=5.0cm,
  xbar, bar width=13pt,
  xmin=0, xmax=290, ymin=-0.7, ymax=4.7,
  xlabel={\scriptsize Liquidity horizon (days)},
  xlabel style={font=\scriptsize}, tick label style={font=\scriptsize},
  xtick={0,50,100,150,200,250},
  ytick={0,1,2,3,4},
  yticklabels={{Cat.\ 5},{Cat.\ 4},{Cat.\ 3},{Cat.\ 2},{Cat.\ 1}},
  nodes near coords, nodes near coords style={font=\tiny\sffamily, color=black!70},
  every node near coord/.append style={/pgf/number format/fixed},
  axis line style={draw=black!55}, clip=false]
\addplot[fill=PaleNavy, draw=FRMNavy!65, line width=0.4pt] coordinates {
(250,0)(120,1)(60,2)(20,3)(10,4)};
\end{axis}
\end{tikzpicture}
\end{center}

\begin{center}
\small
\begin{tabular}{l r >{\raggedright\arraybackslash}p{7.4cm}}
\arrayrulecolor{FRMRule}\toprule
\rowcolor{FRMNavy} \hdr{Category} & \hdr{Horizon} & \hdr{Example asset class}\\
Category 1 & 10 days & Most liquid equities\\
\rowcolor{FRMSoft} Category 2 & 20 days & Investment-grade sovereign credit spreads\\
Category 3 & 60 days & Medium liquidity assets\\
\rowcolor{FRMSoft} Category 4 & 120 days & Non-investment grade corporate credit spreads\\
Category 5 & 250 days & Structured products\\
\bottomrule
\end{tabular}
\end{center}
\figcap{The horizons are not evenly spaced: the jump from category 4 to category 5 is larger
than the whole span from category 1 to category 3. Structured products are assumed to take
twenty-five times as long to exit as liquid equities.}

\subsection{Calculating ES across the horizons}

\begin{center}
\small
\begin{tabularx}{\linewidth}{@{}p{3.0cm} >{\raggedright\arraybackslash}X@{}}
\arrayrulecolor{FRMRule}\toprule
\rowcolor{FRMNavy} \hdr{Approach} & \hdr{Method}\\
\textbf{Old / initial idea} & Use \term{overlapping time periods} for stress testing. Measure
  stressed P\&L changes over multiple trials, scaling up from shorter to longer time periods,
  create a distribution, and calculate ES from it.\\
\rowcolor{FRMSoft}
\textbf{Revised: Internal Models Approach (IMA)} & Use a \term{10-day base horizon} with
  \textbf{nested shocks}: successive shocks to categories, with each ES calculated by holding
  the \textit{more liquid} categories constant. Scaling accounts for nested risk factors, by the
  square root of the time horizon differences.\\
\bottomrule
\end{tabularx}
\end{center}

\begin{center}
\begin{tikzpicture}[every node/.style={font=\tiny\sffamily},
  cell/.style={rectangle, draw=black!40, line width=0.5pt, minimum width=1.55cm,
    minimum height=0.62cm, align=center},
  shk/.style={cell, fill=FRMRed!25, draw=FRMRed!55},
  hld/.style={cell, fill=black!8, draw=black!30}]
  \foreach \c [count=\i from 0] in {1,2,3,4,5}{
    \node[font=\tiny\sffamily, text=black!65] at (\i*1.6+1.9,0.62) {Cat.\ \c};}
  % ES1
  \node[anchor=east, font=\scriptsize\sffamily, text=FRMNavy] at (0.95,0) {$ES_1$};
  \foreach \i in {0,1,2,3,4}{\node[shk] at (\i*1.6+1.9,0) {shock};}
  % ES2
  \node[anchor=east, font=\scriptsize\sffamily, text=FRMNavy] at (0.95,-0.75) {$ES_2$};
  \node[hld] at (1.9,-0.75) {held};
  \foreach \i in {1,2,3,4}{\node[shk] at (\i*1.6+1.9,-0.75) {shock};}
  % ES3
  \node[anchor=east, font=\scriptsize\sffamily, text=FRMNavy] at (0.95,-1.5) {$ES_3$};
  \foreach \i in {0,1}{\node[hld] at (\i*1.6+1.9,-1.5) {held};}
  \foreach \i in {2,3,4}{\node[shk] at (\i*1.6+1.9,-1.5) {shock};}
  % ES4
  \node[anchor=east, font=\scriptsize\sffamily, text=FRMNavy] at (0.95,-2.25) {$ES_4$};
  \foreach \i in {0,1,2}{\node[hld] at (\i*1.6+1.9,-2.25) {held};}
  \foreach \i in {3,4}{\node[shk] at (\i*1.6+1.9,-2.25) {shock};}
  % ES5
  \node[anchor=east, font=\scriptsize\sffamily, text=FRMNavy] at (0.95,-3.0) {$ES_5$};
  \foreach \i in {0,1,2,3}{\node[hld] at (\i*1.6+1.9,-3.0) {held};}
  \node[shk] at (4*1.6+1.9,-3.0) {shock};
  \node[anchor=north, text=black!60, align=center, text width=9.5cm] at (5.1,-3.5)
    {each row shocks one fewer category, always dropping the \textbf{most liquid} one};
\end{tikzpicture}
\end{center}
\figcap{The nesting, laid out. $ES_1$ shocks everything on a 10-day horizon. $ES_2$ holds
category 1 constant and shocks categories 2 to 5. Each successive row freezes one more of the
liquid categories, so only the progressively less liquid exposures keep moving. The formula
below then combines them.}

\begin{mrfmlbox}
\[ ES \;=\; \sqrt{\,ES_1^{2} \;+\; \sum_{j=2}^{5}\left( ES_j
   \sqrt{\frac{LH_j - LH_{j-1}}{10}} \right)^{\!2}} \]
\mrvk{$ES_1$ = the 10-day ES with all categories shocked;\quad $ES_j$ = the ES with categories $j$
through 5 shocked and the more liquid ones held constant;\quad $LH_j$ = the liquidity horizon of
category $j$;\quad the division by 10 rebases each horizon difference onto the 10-day base, and
the square root is the usual time scaling.}{}
\end{mrfmlbox}

\begin{mrkeybox}[title={Until the IMA is approved}]
Banks currently use the \term{Revised Standardized Approach}, which groups assets into buckets
based on risk and then applies a formula using asset values, weights and correlations to produce
a standardized risk measure for each bucket.
\end{mrkeybox}

\subsection{The standardized approach, in three components}

\begin{center}
\begin{tikzpicture}[every node/.style={font=\scriptsize\sffamily},
  top/.style={rectangle, rounded corners=3pt, draw=FRMGold!70, fill=FRMPaper,
    line width=0.9pt, text width=4.8cm, align=center, inner sep=6pt},
  comp/.style={rectangle, rounded corners=3pt, line width=0.8pt,
    text width=3.4cm, align=center, inner sep=6pt, minimum height=1.1cm},
  sub/.style={rectangle, rounded corners=2pt, draw=FRMNavy!45, fill=PaleNavy,
    line width=0.6pt, text width=3.1cm, align=center, inner sep=4pt}]
  \node[top] (cap) at (5.2,2.2) {\textbf{Capital requirement}\\ \textit{sum of three components}};
  \node[comp, draw=FRMNavy!60, fill=PaleNavy] (rs) at (0,0)
    {\textbf{1. Risk sensitivity}\\ seven risk classes};
  \node[comp, draw=FRMGreen!60, fill=PaleGreen] (dr) at (5.2,0)
    {\textbf{2. Default risk charge}\\ jump-to-default};
  \node[comp, draw=FRMRed!60, fill=PaleRed] (rr) at (10.4,0)
    {\textbf{3. Residual risk}\\ e.g.\ exotic options};
  \foreach \t in {rs,dr,rr}{
    \draw[FRMGold, -{Stealth[length=4.5pt]}, line width=0.9pt] (cap) -- (\t);}
  \node[sub] (d) at (0,-1.45) {\textbf{Delta} risk charge};
  \node[sub] (v) at (0,-2.25) {\textbf{Vega} risk charge};
  \node[sub] (c) at (0,-3.05) {\textbf{Curvature} risk charge};
  \draw[black!35, line width=0.6pt] (rs.south) -- (0,-1.13);
\end{tikzpicture}
\end{center}

\begin{center}
\small
\begin{tabularx}{\linewidth}{@{}p{2.5cm} >{\raggedright\arraybackslash}X@{}}
\arrayrulecolor{FRMRule}\toprule
\rowcolor{FRMNavy} \hdr{Charge} & \hdr{What it measures}\\
\textbf{Delta} & The sensitivity of the portfolio to \term{percentage changes in the prices} of
  underlying assets.\\
\rowcolor{FRMSoft}
\textbf{Vega} & The sensitivity of the value of the portfolio to \term{small changes in
  volatility}.\\
\textbf{Curvature} & A capital charge for the bank's \term{gamma} risk exposure.\\
\rowcolor{FRMSoft}
\textbf{Residual} & Risks that \textit{cannot} be handled by the delta/vega/curvature approach,
  such as exotic options.\\
\textbf{Default risk} & \term{Jump-to-default} risk, calculated by multiplying each exposure by
  \textbf{LGD} and a \textbf{default risk weight}, both specified by the Basel Committee.
  Credit \textit{spread} risk, by contrast, is handled through delta/vega/curvature.\\
\bottomrule
\end{tabularx}
\end{center}

\begin{mrkeybox}[title={The seven risk classes}]
General interest rate risk, foreign exchange risk, commodity risk, equity risk, and \textbf{three}
categories of credit spread risk --- corresponding to trading desks.
\end{mrkeybox}

% ---------------------------------------------------------
\lo{MR-18 c}{Explain the FRTB revisions to Basel regulations in the following areas: classification of positions in the trading book compared to the banking book; backtesting; profit and loss attribution; credit risk; and securitizations.}

\subsection{1. Trading book versus banking book}

\begin{center}
\small
\begin{tabularx}{\linewidth}{@{}p{2.4cm} >{\raggedright\arraybackslash}X >{\raggedright\arraybackslash}X@{}}
\arrayrulecolor{FRMRule}\toprule
\rowcolor{FRMNavy} \hdr{} & \hdr{Trading book} & \hdr{Banking book}\\
\textbf{Valuation} & Marked to market & Held to maturity\\
\rowcolor{FRMSoft}
\textbf{Capital} & \textbf{Lower} capital requirements & \textbf{Higher} capital requirements\\
\textbf{Entry test} & Must meet \textit{both}: (a) the bank has the \textbf{ability to trade}
  the asset, and (b) the bank \textbf{physically manages the risk} of the asset on the trading
  desk. Day-to-day price changes must also affect the bank's equity and solvency. & ---\\
\bottomrule
\end{tabularx}
\end{center}

\begin{mrtrapbox}[title={Restrictions on reclassification}]
\begin{itemize}
\item Once assigned to either book, assets \term{cannot be reclassified}, except in
      extraordinary circumstances --- a firm-wide shift in accounting practice, for instance.
\item Any capital benefit from shifting an asset is \term{disallowed}: the original capital
      requirement must be retained.
\end{itemize}
\medskip
The purpose is to prevent \textbf{regulatory arbitrage}, enhance risk management, and improve
transparency. The lower capital charge on the trading book is precisely what creates the
incentive to game the boundary.
\end{mrtrapbox}

\subsection{2. Backtesting}

\begin{mrkeybox}[title={Why stressed ES is not backtested}]
\begin{enumerate}[label=\textbf{\alph*)}]
\item It is more difficult to backtest \term{ES} compared to VaR.
\item It is \term{not possible} to backtest a \textit{stressed} measure at all.
\end{enumerate}
So FRTB backtests a \textbf{VaR} measure instead, even though capital is based on ES.
\end{mrkeybox}

\begin{center}
\small
\begin{tabular}{l >{\raggedright\arraybackslash}p{6.2cm} r}
\arrayrulecolor{FRMRule}\toprule
\rowcolor{FRMNavy} \hdr{Requirement} & \hdr{Detail} & \hdr{Exception limit}\\
Horizon & One day & ---\\
\rowcolor{FRMSoft} Window & Most recent 12 months of data & ---\\
Confidence & 99\% & more than \textbf{12}\\
\rowcolor{FRMSoft} Confidence & 97.5\% & more than \textbf{30}\\
\midrule
Scope & \multicolumn{2}{>{\raggedright\arraybackslash}p{8.6cm}}{\textbf{Each trading desk}
  backtests separately. Both confidence levels must be used.}\\
\bottomrule
\end{tabular}
\end{center}
\figcap{Breaching either limit forces the desk onto the \textbf{standardized approach} for
capital until neither condition continues to exist. Note the pairing: the \textit{lower}
confidence level carries the \textit{higher} exception allowance, because a 97.5\% model is
expected to be breached more often.}

\subsection{3. Profit and loss attribution}

Banks must compare the actual profit or loss in a day with that predicted by their models. Two
measures are calculated.

\begin{mrfmlbox}
\[ \text{Measure 1} \;=\; \frac{\text{mean of } D}{\text{standard deviation of } A}
   \qquad\qquad
   \text{Measure 2} \;=\; \frac{\text{variance of } D}{\text{variance of } A} \]
\mrvk{$D$ = the difference between actual and model profit/loss;\quad $A$ = the actual profit/loss
for a given day. Measure 1 tests for \textit{bias} --- is the model systematically off. Measure
2 tests for \textit{noise} --- how much of the variation the model fails to explain.}{}
\end{mrfmlbox}

\begin{center}
\begin{tikzpicture}[every node/.style={font=\scriptsize\sffamily}]
% measure 1 band
  \draw[black!45, line width=0.7pt] (0,0) -- (11.0,0);
  \fill[PaleGreen, draw=FRMGreen!60, line width=0.7pt] (4.15,-0.32) rectangle (6.85,0.32);
  \foreach \x/\l in {0.6/-30\%, 2.4/-20\%, 4.15/-10\%, 5.5/0, 6.85/+10\%, 8.6/+20\%, 10.4/+30\%}{
    \draw[black!45] (\x,-0.12) -- (\x,0.12);
    \node[anchor=north, text=black!60] at (\x,-0.36) {\l};}
  \node[text=FRMGreen] at (5.5,0) {\textbf{pass}};
  \node[anchor=west, text=FRMRed] at (0.1,0.62) {fail};
  \node[anchor=east, text=FRMRed] at (10.9,0.62) {fail};
  \node[anchor=south, text=black!70] at (5.5,0.62)
    {\textbf{Measure 1} must sit between $\pm 10\%$};
% measure 2 band
  \draw[black!45, line width=0.7pt] (0,-2.7) -- (11.0,-2.7);
  \fill[PaleGreen, draw=FRMGreen!60, line width=0.7pt] (0.6,-3.02) rectangle (4.15,-2.38);
  \foreach \x/\l in {0.6/0, 2.4/10\%, 4.15/20\%, 6.0/30\%, 7.8/40\%, 9.6/50\%}{
    \draw[black!45] (\x,-2.82) -- (\x,-2.58);
    \node[anchor=north, text=black!60] at (\x,-3.06) {\l};}
  \node[text=FRMGreen] at (2.4,-2.7) {\textbf{pass}};
  \node[anchor=west, text=FRMRed] at (4.35,-2.08) {fail};
  \node[anchor=south, text=black!70] at (5.5,-2.08)
    {\textbf{Measure 2} must be under $20\%$};
\end{tikzpicture}
\end{center}
\figcap{Measure 1 is two-sided because bias in either direction is a failure. Measure 2 is
one-sided because there is no such thing as too little unexplained variance. If either ratio
falls outside its band on \textbf{four or more occasions} in a 12-month period, capital must be
computed using the standardized approach.}

\subsection{4. Credit risk}

FRTB provides a modification of the IRC. For instruments dependent on the credit risk of a
particular company, \term{two} types of risk are identified.

\begin{center}
\small
\begin{tabularx}{\linewidth}{@{}p{3.0cm} >{\raggedright\arraybackslash}X@{}}
\arrayrulecolor{FRMRule}\toprule
\rowcolor{FRMNavy} \hdr{Risk type} & \hdr{Treatment}\\
\textbf{a) Credit spread risk} & The risk that a credit asset's spread changes, affecting its
  mark-to-market value. Managed through the \term{expected shortfall} calculation process.\\
\rowcolor{FRMSoft}
\textbf{b) Jump-to-default risk} & The risk of immediate and significant loss due to the default
  of an issuer. Subject to an \term{incremental default risk (IDR) charge}, which applies to
  \textit{all} risk assets including equities, and is calculated using a \textbf{99.9\% VaR with
  a 1-year horizon}.\\
\bottomrule
\end{tabularx}
\end{center}

\begin{mrtrapbox}[title={The IRC replacement assumption}]
Under the old \term{IRC} process, banks assumed a \textbf{constant level of risk} by replacing
deteriorated positions with similar risk assets. For example, an A-rated bond with a three-month
liquidity horizon suffering a credit loss would be replaced with another A-rated bond.

\medskip
\term{FRTB change:} simplifying assumptions like replacement are dropped, in favour of
\textbf{incremental marking to market without assuming replacement}.
\end{mrtrapbox}

\subsection{5. Securitizations}

The \term{Comprehensive Risk Measure (CRM)} was introduced in Basel II.5 to address the risks of
securitized products such as asset-backed securities and collateralized debt obligations.

\begin{center}
\begin{tikzpicture}[every node/.style={font=\scriptsize\sffamily},
  st/.style={rectangle, rounded corners=3pt, line width=0.9pt,
    text width=4.3cm, align=center, inner sep=6pt, minimum height=1.6cm}]
  \node[st, draw=FRMRed!60, fill=PaleRed] (old) at (0,0)
    {\textbf{Basel II.5}\\ banks use their \textit{own models} for CRM, subject to approval};
  \node[st, draw=black!40, fill=FRMSoft] (prob) at (5.4,0)
    {\textbf{The problem}\\ widely varying capital charges across banks for the
     \textit{same} portfolio};
  \node[st, draw=FRMGreen!60, fill=PaleGreen] (new) at (10.8,0)
    {\textbf{FRTB}\\ mandates the \textit{standardized approach} for securitizations};
  \draw[FRMGold, -{Stealth[length=5pt]}, line width=1.1pt] (old) -- (prob);
  \draw[FRMGold, -{Stealth[length=5pt]}, line width=1.1pt] (prob) -- (new);
\end{tikzpicture}
\end{center}
\figcap{The motivation is consistency, not conservatism. Internal models were not judged too
lenient; they were judged too \textit{variable}, so the same portfolio attracted different
capital at different banks.}

\begin{mrkeybox}[title={The pattern across all five areas}]
Read the five revisions together and one theme runs through every one of them: FRTB narrows the
scope for \term{bank discretion}. The trading-book boundary is fixed, reclassification is
barred, desks that fail backtests lose their internal model, desks that fail P\&L attribution
lose their internal model, and securitizations lose internal models outright.
\end{mrkeybox}
